\title{\emph{\zihao{2} 报告正文}}
\maketitle
\thispagestyle{empty}

\ContentDes{（一）立项依据：}
% intentional line break

\NsfcNote{（为什么要开展此项研究，研究的科学技术价值如何）}

\section{研究背景及科学意义}

\subsection{雾计算及移动点对点流媒体对我国信息产业具有重要意义}

随着移动智能设备（尤其是智能手机）硬件性能的不断提升、移动网络的普及，
对于移动智能设备和网络的研究的迅速发展，不少技术也已广泛应用。
一方面，商业化的、成熟的网络应用和服务已经相当普及。
我国的网民数量和移动互联网用户数量已稳居世界首位。
移动网络服务成为现代社会人们日常生活中不可或缺的一部分。
另一方面，移动网络的新应用、新技术、新挑战也层出不穷。
雾计算（Fog Computing）及移动点对点流媒体（Mobile P2P Streaming）
的研究和应用已经成为新的热点、得到了世界各国政府部门、工业界及学术界的极大关注。

目前，移动智能电话的销量已经超过了个人电脑。
思科（Cisco）预测，到2020年，人均拥有的智能网络设备（包括智能手机、平板电脑和机顶盒等）将达到6.58。
如何将各种智能网络设备有效地集成到互联网已经成为了下一代网络技术的重要任务。
雾计算作为有效集成现有互联网和智能网络设备的解决方案，将在未来的网络技术中发挥不可替代的作用。

2016年10月09日，习主席在主持中共中央政治局集体学习时强调：
“加快推进网络信息技术自主创新，朝着建设网络强国目标不懈努力。”
深圳市信息化发展“十三五”规划中也明确指出，
到2020年，深圳市将建成国际一流的信息基础设施，成为“一带一路”的重要信息通信节点和重要的国际信息港。
深圳市已经拥有了移动网络技术从基础研究、应用研发到产业化的完整链条。
牢牢掌握移动网络技术的发展的新热点、新方向，对于进一步巩固和发扬深圳市的产业优势有着深远的战略意义。

\subsection{雾计算及移动点对点流媒体具有挑战性的问题}

互联网中的绝大部分传输数据都是多媒体（尤其是视频）数据，多媒体数据占用网络带宽的比例还在逐年不断增长。
那么，作为网络服务中的重要方面，如何有效地利用现有的网络资源，经济合理地储存和传输大量的多媒体数据，
并将多媒体数据高效地从内容提供者（Content Provider）传输到用户（Client）,则是多媒体网络研究的核心问题。
现有的网络多媒体内容提供者常常利用内容分发网络（CDN: Content distribution network）
或者云计算（Cloud Computing）来支持多媒体服务。
然而，不断增长的多媒体数据和用户需求、更高清晰度的视频标准给数据的储存、传输和处理带来了越来越大的挑战。

对未来的网络，主要挑战如下：
\begin{enumerate}
\item 高质量流媒体服务的数据传输本身需要更高网络带宽的支持。
更高解析度的视频（例如4k）需要更高的数据带宽。
新的多媒体服务（例如虚拟现实技术（Virtual Reality）、增强现实技术（Augmented Reality））
对于带宽的要求也是只增不减。因此，需要研究更有效使用网络资源的技术。

\item 无处不在的移动网络也大大地带动了用户对流媒体服务的需求。
各种各样的智能移动设备将强大的计算和通信能力带到人们的掌中，人们可以有更大的自由度来使用互联网服务。
移动设备上的新型应用技术将极大的改变人们的生活方式，同时也会带来更大的、更多变的用户需求。

\item 移动智能设备本身的移动特性也使得网络结构更加的复杂和多变。
相比较有线网络而言，无线网络带宽依然比较昂贵。
Wi-Fi虽然有较低的可靠性和安全性，却依然被广泛使用。
多媒体服务所需要的高带宽可能需要同时通过有线网络、Wi-Fi和移动网络传输。
有效的服务需要根据实际情况选择最合理的传输方式。
\end{enumerate}

下一代网络技术（有线及无线网络技术）必须解决这三方面因素带来的数据流量对网络的冲击。
传统的内容分发网络和云计算只能通过部署海量昂贵的网络资源来满足用户的需求，
而无限制的在云端部署更多资源并不实际。为了支持未来高质量流媒体（例如高清视频）服务，
必须依靠新的网络条件，有效利用更加普及的智能网络设备。
作为新兴技术，雾计算及移动点对点流媒体解决上述挑战的潜力很大。

雾计算拓展了云计算解决问题的哲学。不同于将资源集中在远端服务器的云计算，
雾计算旨在离用户更近的地方（例如家庭路由器，Wi-Fi热点）向用户提供网络资源和服务。
更短的传输距离不仅可以减少传输的延迟，而且可以减少数据传输量，从而更有效地利用网络资源。
因此，雾计算可以同云计算和无线网络协作，提供更好的用户体验。

移动点对点流媒体技术利用智能手机的资源为其它移动用户提供媒体服务。
在无线网络中，智能手机不仅可以接收媒体数据，也可以分享数据给其它用户。
当多个移动用户收看同一媒体节目时，一个用户可以将自己收到的数据通过Wi-Fi转发给其它用户，
不仅节省了骨干网络的带宽，扩大了信号覆盖面，用户间的分享也提高了传输的可靠性。

\section{总结}

综上所述，本项目拟结合我们已有的多媒体传输网络，基于雾计算和移动点对点流媒体技术，
研究实用化的、可商业部署的、高可靠性的流媒体技术。
研究成果不仅涵盖理论分析和算法设计，而且也会完成能在现有智能手机上使用的软件。

\section{国内外相关工作}

未来流媒体的发展趋势主要有两个新方向。
一方面，视频的清晰度越来越高，未来的高清视频将消耗数倍于目前视频的带宽。
另一方面，新的应用技术层出不穷。
虚拟现实技术和增强现实技术可以向用户提供“浸入式”的多媒体体验，全方位的视角对网络资源的要求也是水涨船高。

目前，往往通过以下技术支持多媒体服务：
\begin{itemize}
\item 内容分发网络（CDN: Content distribution network）：不少大型网络公司拥有分布式的大型服务器，
而这些服务器往往部署在地理上相距甚远的位置，以便能够最大程度的覆盖用户。
这种技术的主要目的是通过分布式地部署一些服务器，使得每一个用户都能接近某个服务器，
从而减少用户和服务器之间传输延迟。

\item 云计算（Cloud Computing）：云计算的核心思想是“把计算作为服务”。
不同于内容分发网络，云计算中服务器的拥有者和使用者通常并不一致。
内容提供商可以根据自己的实际需要，从云计算服务提供商那里“租借”网络资源。
这样一来，资源的规模可以快速的调整，可以有效响应用户的需求，避免资源浪费。

\item 点对点技术（P2P: Peer to peer network）：点对点技术不仅和雾计算技术密切相关，
而且被看作是雾计算技术的前辈。点对点技术的核心是“去中心化”，即任何一个节点都可以向其它节点传输数据，
所以任何一个节点离开网络也不会影响整体网络的功能。雾计算相对于点对点技术的进步之处，
在于避免了刻意强调“去中心化”，各节点之前可以进行更复杂的协调。
同时，雾计算的节点种类更广泛，可以包括手机、路由器和机顶盒等。
\end{itemize}

上述技术用于多媒体服务已被国内外广泛研究，并且已经达到了商业化的水平。现有的视频网站和多媒体网络应用都依赖于上述的技术。

雾计算是近年来兴起的概念，其核心思想是“更靠近用户”。不同于传统的服务器，
雾计算往往利用更小型、更靠近用户的设备（例如智能手机、机顶盒及无线路由器等）来提供网络服务。
最初，为了解决物联网（IoT：Internet of Things）中面临的需求，
这些部署在离用户更近的地方的计算和网络资源能够节省带宽、减小延迟、提高保密性。
在物联网领域，现已有不少国内外的论文和应用产品。目前，国内外开始研究在更广泛的领域使用雾计算技术。
例如，雾计算可以应用于车载网络上，有效报告路面状况。
虽然流媒体服务已成为网络服务的重要组成部分，大部分服务依然是基于“服务器-用户”(Server - Client)的模式，
即用户直接从某个固定的服务器下载多媒体文件。
即使某些流媒体服务利用了部分的用户资源，这种利用也是临时的：靠近用户的网络资源并没有被有效的利用起来。
目前国内外尚没有实用的、可部署的产品使用“雾计算”来提供可靠的流媒体服务。

通过点对点技术分享文件已有几十年的历史。
学术界对此已经有着广泛而深刻的研究，工业界也有成熟的技术（例如：Bit Torrent）。
然而，传统的点对点技术也存在着明显的缺点。点对点技术往往缺乏有效的网络拓扑优化方案，
可能导致过于复杂的拓扑结构。这会大量消耗网络带宽，也会增加流媒体传输的延迟。
同时，在移动设备上（例如智能手机）通过Wi-Fi使用点对点技术传输流媒体尚未有成熟的方案。
由于Wi-Fi传输协议（Protocol）的限制，并非任意两部移动设备都可以利用点对点技术连接。
因此，无法直接将现有的点对点技术移植到移动平台，解决这些问题需要新的设计和算法。

近几年来，我们在流媒体技术的理论上和应用上做了相当的积累。
我们在视频点播（Video-on-Demand）和流媒体直播（Live Streaming）的技术上做了大量的研究工作。
我们发表了多篇相关论文并部署了商业化的视频点播和流媒体直播服务，为本领域的深入研究奠定了良好的基础。
通过对国内外现有研究工作的分析，结合我们已有的基础，
从节省网络资源、支持大规模用户、支持高视频码率、高可靠性等优化目标出发，
提出新的基于雾计算和移动点对点的流媒体传输网络的构造算法，并结合实际的平台进行部署和验证，
最终完成实用化、商业化的系统。

在相关工作的评述中，应该尽量广泛的包含各方面的先进成果：
既有国际上顶级的研究，也得有国内最先进的成果。
\myEmph{建议同时包含自己的中文\cite{21SC_WebSegE}和英文\cite{ChengZMHH10}的代表作，
这些相关的代表作也标明自己在该领域的研究基础。}

{
\bibliographystyle{nsfc.bst}
\bibliography{Cmm}
}
